\section{Introducción}
\label{sec:introduccion}

El objetivo de este trabajo es adquirir el electrocardiograma de un voluntario con un sistema de bajo costo, en reposo y después de una actividad física intensa, y estudiar cómo cambia la variabilidad de su frecuencia cardíaca entre ambos estados. Para eso hay que resolver una cadena completa de procesamiento: entender qué guarda el sistema de adquisición y a qué frecuencia muestrea, acondicionar la señal, detectar cada latido, construir la serie de intervalos entre latidos (tacograma) y estimar su contenido espectral.

El informe se organiza de la siguiente manera. En esta sección se resumen los fundamentos del electrocardiograma y sus derivaciones, los formatos en que se almacena y las bases fisiológicas de la variabilidad de la frecuencia cardíaca. La sección~\ref{sec:metodos} describe el protocolo, el sistema de adquisición y cada etapa del procesamiento. La sección~\ref{sec:resultados} presenta los resultados del acondicionamiento, la detección de latidos, los tacogramas y los periodogramas, que se discuten frente a la teoría en la sección~\ref{sec:discusion}. Las conclusiones se resumen en la sección~\ref{sec:conclusiones}. Todo el procesamiento está implementado en Python; el programa \texttt{main.py} del repositorio del grupo lo corre de punta a punta y genera las figuras y tablas de este informe.

\subsection{Electrocardiograma y derivaciones}
\label{sec:derivaciones}

El electrocardiograma (ECG) registra en la superficie del cuerpo las diferencias de potencial que genera la actividad eléctrica del corazón. Cada derivación es la diferencia de potencial entre un electrodo explorador y una referencia, que puede ser otro electrodo o una combinación de varios, y muestra la proyección del vector cardíaco sobre una dirección particular \parencite{malmivuo1995}. Por eso la morfología de las ondas depende de dónde se coloquen los electrodos, y un trazado solo puede compararse con otro, o con los criterios diagnósticos, si se registró con la misma ubicación.

\subsubsection{Electrodos en los miembros}
\label{sec:electrodos-miembros}
% Responde: ítem 3b

En el ECG clínico de reposo se colocan electrodos en ambos brazos (RA y LA) y en ambas piernas (LL y RL). La razón principal es eléctrica. Los miembros se comportan como conductores largos y delgados conectados al tronco por un extremo, por los que casi no circula corriente de origen cardíaco, de modo que el potencial en la muñeca es prácticamente el mismo que en la unión del brazo con el tórax. Medir en la muñeca o en el tobillo equivale entonces a medir en el hombro o en la cadera. Esos tres puntos quedan lejos del corazón y a distancias parecidas de él, y forman el triángulo de Einthoven, sobre el que se definen las derivaciones bipolares del plano frontal \parencite{malmivuo1995}.

La segunda razón es la estandarización. Los valores normales de amplitud y de eje eléctrico, y los criterios diagnósticos que se usan en la práctica, se establecieron con los electrodos en esas posiciones. Si se los desplaza al torso, como en la configuración de Mason-Likar que se usa en las pruebas de esfuerzo, cambian las amplitudes y el eje, y el trazado deja de ser equivalente a uno estándar \parencite{kligfield2007}. Los electrodos de los miembros son, además, la base de las demás derivaciones: RA, LA y LL promediados forman el terminal central de Wilson, que es la referencia de las derivaciones precordiales, y combinados de a dos dan la referencia de las derivaciones aumentadas. El electrodo de la pierna derecha (RL) no forma parte de ninguna derivación. Sirve de referencia del amplificador o, en los equipos con circuito de pierna derecha activa, recibe una señal que reduce la tensión de modo común y con ella la interferencia de línea.

En este trabajo usamos otra configuración: los tres electrodos en el torso, uno debajo de cada clavícula y el tercero en el abdomen, por encima de la cadera derecha. La placa mide la diferencia entre el electrodo izquierdo y el derecho, una aproximación a la derivación I, y conecta el tercero al circuito de pierna derecha activa. Cerca del corazón y lejos de los músculos de los brazos, la señal tiene más amplitud y menos contaminación muscular \parencite{sparkfun_guia}. Esto alcanza para detectar los complejos QRS y medir intervalos RR, pero el registro no reproduce ninguna derivación estándar y su morfología no se puede evaluar con criterios clínicos.

\subsubsection{Otras derivaciones}
\label{sec:otras-derivaciones}
% Responde: ítem 3c

El estudio estándar de reposo registra doce derivaciones. En el plano frontal están las tres bipolares de Einthoven (I, II y III, que miden LA$-$RA, LL$-$RA y LL$-$LA) y las tres aumentadas de Goldberger (aVR, aVL y aVF), que miden el potencial de cada miembro respecto del promedio de los otros dos. Juntas forman el sistema hexaxial: seis direcciones separadas \ang{30} que cubren el plano frontal. En el plano horizontal están las seis precordiales, V1 a V6, que miden el potencial de puntos del tórax ubicados entre el borde derecho del esternón y la línea axilar media izquierda, respecto del terminal de Wilson \parencite{kligfield2007,malmivuo1995}.

Fuera del estudio estándar se agregan derivaciones para situaciones particulares: precordiales derechas (V3R a V6R) para evaluar el ventrículo derecho, posteriores (V7 a V9) para la pared posterior del ventrículo izquierdo, el sistema ortogonal de Frank (X, Y, Z) que se usa en vectorcardiografía, y derivaciones bipolares de tórax modificadas para Holter y monitoreo continuo, del mismo tipo que la que usamos en este experimento.

Con doce vistas distribuidas en dos planos casi perpendiculares, el profesional puede reconstruir la dirección de la activación eléctrica en el espacio y ubicar dónde está una anomalía. El eje eléctrico medio del QRS en el plano frontal orienta hacia hipertrofias o bloqueos fasciculares cuando se desvía. Los cambios del segmento ST y las ondas Q patológicas se localizan según el grupo de derivaciones donde aparecen: cara inferior en II, III y aVF; cara lateral en I, aVL, V5 y V6; región anterior y septal en V1 a V4. A partir de esa localización se infiere qué arteria coronaria está comprometida. Los bloqueos de rama se reconocen por la forma del QRS en V1 y V6, y la hipertrofia ventricular, por criterios que comparan amplitudes entre derivaciones precordiales \parencite{wagner2014}. Una sola derivación, en cambio, permite medir frecuencia y ritmo, que es lo que necesita el análisis de variabilidad de este trabajo, pero no localizar nada.

\subsection{Formatos de almacenamiento de ECG}
\label{sec:formatos}
% Responde: ítem 1 (segunda parte)

Los ECG digitales se guardan en formatos muy distintos según su uso, y lo que más los diferencia es cuánta información acompaña a las muestras. \textcite{bond2011} revisaron los formatos existentes y distinguieron estándares abiertos, formatos propietarios de los fabricantes y formatos de investigación. Entre los estándares clínicos, los más difundidos son:

\begin{itemize}
  \item \textbf{SCP-ECG} (norma europea EN~1064, adoptada luego como ISO~11073-91064): formato binario y compacto para ECG de reposo, organizado en secciones con los datos del paciente y del equipo, la señal (con compresión opcional) y las mediciones e interpretación automáticas.
  \item \textbf{HL7 aECG}: formato XML que la FDA de Estados Unidos adoptó para recibir los ECG de los ensayos clínicos.
  \item \textbf{DICOM Waveform}: extensión del estándar de imágenes médicas que permite archivar el ECG junto con los demás estudios del paciente en los sistemas de almacenamiento de imágenes (PACS).
  \item \textbf{ISHNE}: formato binario pensado para registros Holter de larga duración.
\end{itemize}

En investigación es habitual el formato \textbf{WFDB} de PhysioNet, que separa un encabezado en texto (frecuencia de muestreo, ganancia, resolución del conversor y nombre de cada derivación), un archivo binario con las muestras y archivos de anotaciones con la posición y el tipo de cada latido \parencite{goldberger2000}. También se usa \textbf{EDF}, un formato general para señales fisiológicas multicanal cuyo encabezado guarda, para cada canal, la frecuencia de muestreo y los rangos físico y digital necesarios para convertir las muestras a unidades; su versión EDF+ agrega anotaciones \parencite{kemp2003}. Muchos equipos exportan además en formatos propios del fabricante, en general XML, o directamente en PDF o como imagen del trazado. Estos últimos sirven para leer el ECG, pero no para procesar la señal.

Frente a estos formatos, el archivo CSV de este experimento (sección~\ref{sec:formato}) es el más simple posible: no guarda la frecuencia de muestreo, las unidades, la ganancia, la derivación ni datos del sujeto. Toda esa información tuvimos que obtenerla por fuera del archivo, a partir del firmware y del código de la interfaz.

\subsection{Variabilidad de la frecuencia cardíaca}
\label{sec:hrv}
% Base teórica para los ítems 6, 7 y 8

El corazón late con un ritmo propio, marcado por las células marcapasos del nodo sinusal, pero ese ritmo está modulado en forma continua por el sistema nervioso autónomo. La rama parasimpática, a través del nervio vago, reduce la frecuencia cardíaca y actúa muy rápido: su efecto aparece y desaparece en menos de un segundo, de modo que puede cambiar la duración de un latido al siguiente. La rama simpática la aumenta, pero con una dinámica más lenta, de varios segundos. Por eso el intervalo entre latidos no es constante ni siquiera en reposo, y sus fluctuaciones, la variabilidad de la frecuencia cardíaca (VFC), reflejan la actividad de ambas ramas y de los reflejos que las controlan \parencite{taskforce1996,shaffer2017}.

La serie de intervalos entre picos R consecutivos (intervalos RR) en función del tiempo se llama \textbf{tacograma}. Sobre ella se calculan índices en el dominio del tiempo y de la frecuencia. Entre los primeros, la \textcite{taskforce1996} define el desvío estándar de los intervalos (SDNN), que mide la variabilidad total; la raíz cuadrática media de las diferencias entre intervalos sucesivos (RMSSD) y el porcentaje de diferencias sucesivas mayores a \SI{50}{\milli\second} (pNN50). Estos dos últimos miden la variabilidad de un latido al siguiente, que solo la rama vagal es capaz de producir, y por eso se usan como índices del tono vagal \parencite{shaffer2017}.

En el dominio de la frecuencia, el espectro del tacograma se divide en bandas \parencite{taskforce1996}:
\begin{itemize}
  \item \textbf{HF} (alta frecuencia, \SIrange{0.15}{0.4}{\hertz}): corresponde a la arritmia sinusal respiratoria, el aumento de la frecuencia cardíaca durante la inspiración y su disminución durante la espiración. Está mediada casi por completo por el vago, y la frecuencia de su pico coincide con la frecuencia respiratoria.
  \item \textbf{LF} (baja frecuencia, \SIrange{0.04}{0.15}{\hertz}): incluye las oscilaciones de unos \SI{10}{\second} de período ligadas al barorreflejo. Recibe aportes de las dos ramas, por lo que no es un índice simpático puro.
  \item \textbf{VLF} (muy baja frecuencia, por debajo de \SI{0.04}{\hertz}): mecanismos más lentos (termorregulación, sistema renina-angiotensina, tendencias). No se puede estimar en registros de pocos minutos.
\end{itemize}
El cociente LF/HF se usó durante mucho tiempo como medida del balance simpático-vagal. Esa interpretación es discutida: como la banda LF también tiene aporte vagal y ninguna de las dos bandas es independiente de la frecuencia cardíaca, el cociente no mide directamente el balance entre ambas ramas \parencite{billman2013}. Por eso conviene interpretarlo junto con las potencias absolutas y las unidades normalizadas.

La VFC depende además de la propia frecuencia cardíaca. Cuando la frecuencia aumenta, los intervalos RR se acortan y la misma modulación autonómica produce variaciones absolutas menores, por lo que todos los índices de variabilidad tienden a bajar \parencite{sacha2014}.

\paragraph{Ejercicio y recuperación.}
Al comenzar el ejercicio, la frecuencia cardíaca sube primero por la retirada del tono vagal y después, a intensidades mayores, por la activación simpática y las catecolaminas circulantes. Durante un esfuerzo intenso el control vagal es mínimo y la VFC se reduce mucho, sobre todo en la banda HF. Al terminar el ejercicio el proceso se invierte, pero con dos velocidades: la reactivación vagal es rápida y explica la caída de la frecuencia cardíaca en el primer minuto de recuperación, mientras que la retirada simpática y la eliminación de las catecolaminas son más lentas y mantienen la frecuencia elevada durante varios minutos \parencite{imai1994,michael2017}. Por eso se espera que, inmediatamente después del ejercicio, la frecuencia cardíaca esté alta y la VFC deprimida, y que la variabilidad, en particular la de la banda HF, se recupere en forma progresiva mientras la frecuencia baja. La respiración, que sigue agitada al principio de la recuperación, desplaza además el pico HF hacia frecuencias más altas.